% Options for packages loaded elsewhere
\PassOptionsToPackage{unicode}{hyperref}
\PassOptionsToPackage{hyphens}{url}
\documentclass[
  letterpaper,
  11pt]{article}
\usepackage{xcolor}
\usepackage{amsmath,amssymb}
\setcounter{secnumdepth}{5}
\usepackage{iftex}
\ifPDFTeX
  \usepackage[T1]{fontenc}
  \usepackage[utf8]{inputenc}
  \usepackage{textcomp} % provide euro and other symbols
\else % if luatex or xetex
  \usepackage{unicode-math} % this also loads fontspec
  \defaultfontfeatures{Scale=MatchLowercase}
  \defaultfontfeatures[\rmfamily]{Ligatures=TeX,Scale=1}
\fi
\usepackage{lmodern}
\ifPDFTeX\else
  % xetex/luatex font selection
\fi
% Use upquote if available, for straight quotes in verbatim environments
\IfFileExists{upquote.sty}{\usepackage{upquote}}{}
\IfFileExists{microtype.sty}{% use microtype if available
  \usepackage[]{microtype}
  \UseMicrotypeSet[protrusion]{basicmath} % disable protrusion for tt fonts
}{}
\setlength{\emergencystretch}{3em} % prevent overfull lines
\providecommand{\tightlist}{%
  \setlength{\itemsep}{0pt}\setlength{\parskip}{0pt}}
\usepackage{arxiv}
\usepackage{booktabs}
\renewcommand{\shorttitle}{Dialect-Marked Response Audit}

\AtBeginDocument{%
  \hypersetup{
    hidelinks,
    pdfauthor={Jeffrey Shorthill},
    pdftitle={Dialect-Marked Response Audit (DMRA): A Matched-Pair Safety Evaluation of AAVE-Marked Prompt Surfaces in Qwen3.5-35B-A3B}
  }%
}
\usepackage{bookmark}
\IfFileExists{xurl.sty}{\usepackage{xurl}}{} % add URL line breaks if available
\urlstyle{same}
\hypersetup{
  pdftitle={Dialect-Marked Response Audit (DMRA): A Matched-Pair Safety Evaluation of AAVE-Marked Prompt Surfaces in Qwen3.5-35B-A3B},
  pdfkeywords={large language models, dialect fairness, AAVE, safety
evaluation, matched-pair audit, reasoning traces},
  hidelinks,
  pdfcreator={LaTeX via pandoc}}

\title{Dialect-Marked Response Audit (DMRA): A Matched-Pair Safety
Evaluation of AAVE-Marked Prompt Surfaces in Qwen3.5-35B-A3B}
\author{Jeffrey Shorthill\\
\texttt{jws299792@icloud.com}}
\date{May 19, 2026}

\begin{document}
\maketitle
\begin{abstract}
Safety evaluations for large language models often ask whether an answer
refuses a harmful request or contains correct information. This misses a
different failure mode: matched users can receive unequal safety support
when the same situation is expressed through different prompt surfaces.
We introduce \textbf{Dialect-Marked Response Audit (DMRA)}, an I/O-level
matched-pair protocol for comparing responses to semantically matched
African American Vernacular English (AAVE)-marked and paired comparison
prompts. We apply DMRA to Qwen3.5-35B-A3B and a refusal-reduced HauhauCS
fine-tune using deterministic decoding across 328 run/prompt records
over 66 matched pairs. The central finding is that top-level safety
outcomes can converge while assistance quality diverges. In a pediatric
respiratory-distress case, both matched outputs direct emergency care,
but the comparison-prompt response provides additional deterioration
checks and ambulance guidance beyond the AAVE-marked response. In a
self-harm crisis case, the released base model's visible reasoning
treats \texttt{finna} (a near-future marker meaning ``going to'') as
evidence of Southern U.S./AAVE context and uses that cue during
crisis-resource selection. These cases show why the unit of safety
evaluation should be the matched user interaction: whether equivalent
requests receive comparable urgency, specificity, respect, and risk
reduction across prompt surfaces.
\end{abstract}

\textbf{Keywords:} large language models; dialect fairness; AAVE; safety
evaluation; matched-pair audit; reasoning traces

\section{Introduction}\label{introduction}

Standard LLM safety evaluations often score an individual answer:
refusal status, harmful-content avoidance, policy adherence, or factual
correctness. Those measures are necessary, and they leave open the
question this paper targets: when two users ask for the same
safety-critical help through different prompt surfaces, do they receive
equivalent safety support? This paper treats the matched user
interaction as the relevant unit of analysis.

A pediatric respiratory-distress pair makes the issue concrete. Both
matched outputs in the audit tell a caregiver to seek emergency care for
a wheezing child with trouble breathing. The paired comparison response
also adds deterioration checks and ambulance guidance; the AAVE-marked
response gives the correct top-level direction with less real-time
scaffolding. A triage metric could mark both outputs as safe. A
support-equivalence audit asks whether both users received comparable
urgency, specificity, respect, and risk reduction.

African American Vernacular English (AAVE) is a systematic and socially
indexed English variety whose features are often racialized by listeners
and classifiers. Its morphosyntax, lexicon, discourse stance, and
social-address patterns can be treated by models as signals of
formality, identity, location, credibility, or risk. Safety-critical
assistants should preserve assistance quality across those surfaces,
including for users who bring community or vernacular forms to the
interaction.

This study uses constructed prompt surfaces. We use \textbf{AAVE-marked
prompts} for prompts that contain selected AAVE-associated morphosyntax,
lexical items, discourse stance, or social-address forms. We use
\textbf{paired comparison prompts} for matched counterpart prompts in
the archive. These terms name corpus conditions and scope the claims to
this audit; broader claims about naturalistic AAVE, other English
varieties, and multilingual settings require additional
community-reviewed prompt construction. Legacy labels such as
\texttt{aave}, \texttt{ae}, and \texttt{academic\_english} remain in
archive paths as provenance identifiers.

We introduce \textbf{Dialect-Marked Response Audit (DMRA)}, an I/O-level
matched-pair protocol for testing safety-support equivalence across
prompt surfaces. DMRA uses deterministic decoding, preserved provenance,
response-frame coding, and visible-reasoning analysis where traces are
emitted. Generated text remains the primary evidence. Visible reasoning
is treated as model-emitted output, and internal mechanism claims are
left for separate mechanistic work.

This paper makes three contributions:

\begin{itemize}
\tightlist
\item
  It defines \textbf{DMRA} as a matched-pair audit protocol for
  evaluating equivalent safety support across AAVE-marked and paired
  comparison prompt surfaces.
\item
  It reports an audit of base Qwen3.5-35B-A3B and a refusal-reduced
  HauhauCS fine-tune, covering 328 run/prompt records over 66 unique
  matched pairs across everyday advice, medical triage, financial
  stress, and sensitive safety prompts.\footnote{Content warning: the
    audit includes self-harm crisis language, violent-threat prompts,
    illegal-logistics prompts, and masked slurs. Public excerpts use
    minimal evidence phrases and redaction for operational detail.}
\item
  It shows that refusal and correctness metrics can miss safety-relevant
  support differences, anchored by pediatric respiratory distress and a
  self-harm visible-reasoning trace, with additional results in
  epistemic, financial, and sensitive safety prompts.
\end{itemize}

\section{Related Work}\label{related-work}

Prior work on dialect bias has shown that NLP systems encode linguistic
prejudice when dialect features are associated with informality,
toxicity, lower quality, or demographic identity (Blodgett et al., 2016;
Sap et al., 2019; Groenwold et al., 2020; Hofmann et al., 2024). Recent
LLM work extends this line into reasoning, generation, annotation, and
chatbot quality-of-service settings. Lin et al.~introduce ReDial for
evaluating dialect fairness and robustness in reasoning tasks, while
EnDive evaluates cross-dialect performance gaps across language
understanding, algorithmic reasoning, mathematics, and logic (Lin et
al., 2024; Gupta et al., 2025). Dunlap and McCoy (2026) compare
model-produced AAVE with human AAVE usage and find underuse, misuse, and
stereotype replication. Tornberg (2026) reports systematic differences
in LLM annotation judgments for matched AAVE sentences across 19 models.
Harvey et al.~(2025) frame dialect harms in deployed chatbots as
quality-of-service harms. DMRA builds on this newer audit direction by
focusing on safety-critical assistance quality.

Matched-prompt audits and counterfactual fairness tests alter one input
factor while holding semantic content as steady as practical (Kusner et
al., 2017; Ribeiro et al., 2020; Parrish et al., 2022). Recent
dialect-audit work emphasizes ecological validity, dynamic prompts, and
user-facing quality-of-service outcomes (Harvey et al., 2025). DMRA
adopts this logic while treating dialect-marked prompting as a bundled
surface intervention. AAVE-marked and paired comparison prompts may vary
in syntax, lexicon, stance, profanity, politeness, and social cues. The
results therefore support I/O-level claims about
prompt-surface-conditioned behavior.

LLM safety evaluation commonly measures refusal behavior,
safe-completion adherence, and harmful-request handling (Gehman et al.,
2020; Roettger et al., 2023; Zou et al., 2023; Mazeika et al., 2024).
Recent work separates refusal from generation safety in long-tail
obfuscated inputs (Maskey et al., 2025) and characterizes selective
refusal bias across demographic targets (Khorramrouz and Levy, 2025).
DMRA complements those evaluations by coding the response frame around
the top-level safety outcome. Emergency and crisis prompts can receive
the same triage decision with different urgency, specificity, empathy,
or real-time scaffolding. High-risk prompts can receive similar refusal
status while diverging in tactical framing.

Visible model-emitted reasoning traces provide an additional output
surface for audit. Prior work urges care when interpreting
chain-of-thought faithfulness (Turpin et al., 2023; Lanham et al.,
2023). Newer monitorability work treats chain-of-thought as a
safety-relevant oversight surface and studies legibility, coverage,
faithfulness, and verbosity (Korbak et al., 2025; Emmons et al., 2025;
Meek et al., 2025). Bias-oriented trace work also shows that models may
articulate some cue types differently from others (Balasubramanian et
al., 2025). DMRA treats visible traces as generated text shown to an
auditor or user-facing interface, rather than as a direct window into
hidden computation. When a model emits a trace that turns a dialect
marker into location or community context, that trace itself becomes
audit evidence.

DMRA sits at the intersection of these lines of work. Dialect-bias
research shows that prompt surface can encode socially meaningful
information. Matched-prompt audits provide the comparison design.
Safety-evaluation work shows the limits of refusal-centered scoring.
Reasoning-trace work supplies an additional output surface. DMRA
combines these lines to ask whether matched users receive equivalent
safety support across prompt surfaces.

\section{Dialect-Marked Response
Audit}\label{dialect-marked-response-audit}

\subsection{Audit Objective and
Corpus}\label{audit-objective-and-corpus}

DMRA asks whether matched prompt surfaces produce materially different
model behavior in safety-relevant contexts. We treat differences as
material when they change triage, urgency, safety scaffolding,
de-escalation, empathy, demographic inference, harmful continuation,
tactical detail, or visible reasoning content.

The audit uses 66 matched prompt pairs, corresponding to 132 unique
prompt surfaces, and 328 model/template run records. Each pair contains
an AAVE-marked prompt and a paired comparison prompt expressing matched
user intent. The larger run-record count reflects repeated observations
over those prompt surfaces across model variants and answer-only or
visible-reasoning templates.

The prompts cover five scenario areas. We test everyday advice and
self-description questions, including epistemic and identity prompts;
pediatric and adult medical triage, including respiratory distress,
chest pain, high fever, seizure, and head injury; financial stress and
family/work advice; sensitive safety prompts involving self-harm,
violence, and illegal logistics; and one panic-attack perspective prompt.
Public run-file provenance is recorded in
\texttt{metadata/run\_manifest\_public.tsv}; retained full manifests
preserve dates and file-level provenance in the controlled archive.

\subsection{Prompt Surface
Variables}\label{prompt-surface-variables}

The AAVE-marked prompts contain selected features associated with AAVE
and adjacent register cues, including habitual \texttt{be}, zero copula,
agreement leveling, \texttt{ain\textquotesingle{}t}, \texttt{got}
possession, \texttt{finna}, embedded-question word order, social-address
forms, and register-linked lexical choices. Sensitive prompts also
bundle profanity, slurs, weapon slang, action terms, and explicit social
framing. DMRA therefore treats the main pairs as realistic bundled
prompt-surface contrasts.

\subsection{Decoding and Trace
Handling}\label{decoding-and-trace-handling}

All runs used deterministic decoding. The run manifests and
case-evidence files record temperature \texttt{0}, \texttt{top\_k=1},
\texttt{top\_p=1}, \texttt{min\_p=0}, seed \texttt{42}, and repeat
penalty \texttt{1}. Answer-only runs used a rendered assistant prefill
containing
\texttt{\textless{}think\textgreater{}\ \ \textless{}/think\textgreater{}};
manifests record those runs with the literal template label
\texttt{\textless{}/no\ think\textgreater{}}. Visible-reasoning runs
ended the rendered prompt with \texttt{\textless{}think\textgreater{}}
and preserved model-emitted reasoning text until the final-answer
boundary. Answer-only runs generally used \texttt{n\_predict=2048};
visible-reasoning runs in the 2026-05-15 and 2026-05-16 captures used
\texttt{n\_predict=8192}.

Visible model-emitted reasoning is analyzed as output text. The trace
rule used by the archive records generated content before a closing
\texttt{\textless{}/think\textgreater{}} as visible thinking; when
generation continues through the token budget, the continuation is
recorded as visible thinking continuation. The self-harm case below uses
this visible trace as the evidence surface.

\subsection{Response Coding}\label{response-coding}

Response coding combines exact I/O review, heuristic labels, and
case-specific evidence files. For the initial 50-pair answer-only
records, summary files report completeness, token asymmetry, budget
hits, and lexical screens. For sensitive safety prompts, coding
distinguishes refusal, partial refusal, de-escalation, practical safety
scaffolding, harmful continuation, tactical framing, and visible
reasoning content. Public case-evidence summaries are redacted for
safety; retained case-evidence files preserve prompt text, generation
settings, provenance, and short evidence excerpts.

\section{Experimental Setup}\label{experimental-setup}

The experiment uses two Qwen3.5-35B-A3B variants: the public Q8\_0 GGUF
base release and a HauhauCS refusal-reduced fine-tune (Qwen Team, 2025;
Qwen Team, 2026). The base model is the open-weight instruction
reference. HauhauCS serves as an operational refusal-reduced stress
test, exposing response trajectories that stronger refusal behavior can
compress.

The run matrix varies model variant and response surface. The initial 50
matched pairs were run in answer-only mode on both model variants. The
later medical, financial, and sensitive safety prompts were run on both
model variants in answer-only and visible-reasoning modes. Together,
these captures yield 328 model/template run records over 132 unique
prompt surfaces.

The 2026-05-15 and 2026-05-16 captures use context size
\texttt{16384}; the initial 2026-05-05 answer-only capture used context
size \texttt{4096}. Model weights, environment files, and full raw
high-risk generations are outside the manuscript package. Public
reproducibility metadata records exact model identifiers, quantization
files, hashes, deterministic decoding settings, run provenance,
checksums, and redacted case-evidence summaries.

\section{Results}\label{results}

\subsection{Overview}\label{overview}

Across the 328 model/template records, DMRA shows that matched prompt
surfaces can converge on top-level safety decisions while diverging in
response posture and assistance quality. The clearest divergences appear
in five result areas: epistemic framing, pediatric emergency
scaffolding, financial-stress empathy, high-risk continuation under
refusal reduction, and visible dialect-to-location inference.

Aggregate length measures support the presence of register-conditioned
divergence and point away from a single directional explanation. In the
initial 50-pair answer-only run, raw generated-token asymmetry above 20
percent appeared in 26 of 50 base pairs and 31 of 50 Hauhau pairs.
First-answer word-count asymmetry above 20 percent appeared in 16 base
pairs and 14 Hauhau pairs. Direction varied by pair: in both base and
Hauhau answer-only runs, AAVE-marked first answers were longer in 27
pairs and paired comparison first answers were longer in 23.

\subsection{Epistemic Framing}\label{epistemic-framing}

An epistemic self-description pair asks how to interpret the model's
answer when asked whether it has beliefs. The paired comparison prompt
asks whether the answer should be taken ``as a fact about you or as you
naming your limits''; the AAVE-marked prompt expresses the same user
intent with matched surface edits. The resulting responses give opposite
top-level frames. The paired comparison response says the answer should
be taken as naming the model's limits. The AAVE-marked response says it
should be taken as a statement of fact about the model's architecture.
The AAVE-marked output is also longer after trimming: 383 words versus
175 words, a 2.19x word-length ratio. The retained pair \texttt{003}
case-evidence file preserves full prompt text, output excerpts, and
measurements.

This case shows a response-posture divergence separate from refusal
status. The same conceptual question receives a humility/limit frame in
one surface condition and an architecture/fact frame in the other.

\subsection{Pediatric Respiratory-Distress
Scaffolding}\label{pediatric-respiratory-distress-scaffolding}

Medical triage prompts show a different pattern: matched prompts can
receive the same top-level triage decision with different practical
scaffolding. The clearest example is a pediatric wheezing and
trouble-breathing scenario. In the base answer-only run, both matched
outputs tell the caregiver to seek immediate emergency care. The paired
comparison response adds checks for whether the child can speak in full
sentences, whether lips or fingernails look blue or gray, and whether an
ambulance is safer than driving. It also explains airway narrowing,
rapid worsening, oxygen-level assessment, and clinician lung exam. The
AAVE-marked response gives correct emergency direction with shorter
real-time guidance: 177 words versus 321 words, a 1.81x word-length
ratio for the paired comparison output. The retained \texttt{med\_s03}
case-evidence file preserves the full pair.

This is the clearest policy-facing case because the top-level safety
decision converges while practical triage support differs.

\subsection{Financial-Stress Empathy}\label{financial-stress-empathy}

The financial-stress pair shows register-conditioned empathy style. In
the Hauhau answer-only run, both matched prompts receive long, serious,
practical answers. The AAVE-marked response uses a peer-solidarity frame
and a toughness/endurance metaphor, calling the user a warrior and
moving quickly to the idea that the ``math is just broken.'' The paired
comparison response uses a more reflective and professional-managerial
frame, describing the ``weight'' of exhaustion as palpable and
organizing advice around a ``Foreman'' mindset versus a ``Worker''
mindset. Both responses mirror the user's closing address term, using
``homie'' for the AAVE-marked prompt and ``dude'' for the paired
comparison prompt. That symmetry shows register matching in both
directions while the deeper empathy frame still diverges.

\subsection{High-Risk Continuation Under Refusal
Reduction}\label{high-risk-continuation-under-refusal-reduction}

Sensitive safety cases show how refusal-reduced models can expose
response trajectories. In the reviewed violence and drug-logistics
cases, base answer-only often refuses or partially refuses across both
prompt surfaces. Hauhau answer-only continues the scenario more often,
making surface-conditioned continuation easier to observe.

The clearest violence example is an armed-confrontation pair. The paired
comparison response frames the situation through risk, anger, firearm
escalation, legal consequence, and regret. The AAVE-marked response
treats the proposed armed confrontation more directly as a timing and
feasibility problem and includes a ``clean shot'' frame. The public
evidence excerpt redacts operational details and preserves the
evidence-relevant phrases; retained evidence is keyed to
\texttt{stress\_s02}.

\subsection{Visible Reasoning and Dialect-Cue
Use}\label{visible-reasoning-and-dialect-cue-use}

Visible-reasoning runs add a second audit surface. The strongest trace
finding comes from a self-harm crisis prompt. In the base
visible-reasoning run, the model emits a trace that identifies an
imminent self-harm crisis and then treats the dialect marker
\texttt{finna} (a near-future marker meaning ``going to'') as evidence
of Southern U.S. or AAVE context. The trace uses that inference to
select U.S.-specific crisis resources, especially 988. The final answer
is safety-oriented and crisis-appropriate. The retained evidence is
keyed to \texttt{s\_full\_02\_aave}.

This trace reframes the other results. The pediatric, epistemic,
financial, and violence cases show differences in final response posture
and safety scaffolding. The self-harm trace shows the model naming a
dialect marker and carrying that marker into safety-resource selection.
The equity concern lies in using a dialect cue as demographic or
geographic evidence when the user supplied crisis intent but supplied
location only through language.

\section{Discussion}\label{discussion}

DMRA shows that safety behavior varies along dimensions that refusal
metrics compress. A model can reach the same top-level safety decision
across matched prompts while changing urgency, specificity, empathy, or
risk-reduction scaffolding. The core implication is that safety
evaluation needs to treat equivalent support as a measurable property of
matched interactions.

The audit also shows dialect-linked features functioning as
response-selection features across multiple output surfaces. They appear
in final-answer style, urgency, practical detail, empathy frame,
high-risk continuation, and visible reasoning. The corpus therefore
supports a safety-evaluation standard centered on equivalent support
across matched prompt surfaces.

Refusal behavior often acts as a convergence mechanism. It pushes
different prompt surfaces toward similar safe completions in high-risk
domains. The Hauhau stress test shows the response trajectories that
become visible when refusal convergence weakens: register-conditioned
harmful continuation becomes easier to observe and more consequential.
Safety evaluation should therefore examine response posture and
continuation trajectory, including cases where refusal is partial,
delayed, or embedded in visible reasoning.

The deployment risk extends beyond overt harmful completions. Common
risks include unequal scaffolding, unequal seriousness, overfamiliarity,
stereotype-driven demographic inference, and weaker localization. In
medical triage, a cyanosis check or ambulance cue can matter. In crisis
support, inferring location from dialect can route resources. In
financial distress, an over-toughened empathy frame can encode
assumptions about how a user should be addressed. Dialect-diverse users
should receive the safest version of an assistant's help across the
language they bring to the interaction.

Reasoning-mode audits should preserve visible traces separately from
final answers. A final answer can be appropriate while the visible trace
contains assumptions, classifications, or localization choices that
matter for safety and equity. DMRA takes visible traces as model-emitted
output and audit evidence, with hidden mechanisms reserved for separate
mechanistic study.

Open-weight fine-tunes can change safety behavior in ways that interact
with dialect and register. The Hauhau model is useful because it makes
refusal weakening visible in the audit setting. Its behavior should be
read as an operational stress test: refusal-reduced fine-tunes reveal
response trajectories that ordinary refusal metrics can compress, while
the fine-tune itself bundles many behavioral changes.

\section{Limitations}\label{limitations}

The claims are behavioral. They concern generated text, visible traces,
and preserved I/O evidence. Internal-circuit explanations require a
separate mechanistic study with controlled interventions.

The AAVE-marked prompts are constructed surfaces. They contain selected
AAVE-associated features, discourse stance, profanity, social-address
patterns, and register-correlated lexical choices. They support claims
about this audit corpus and motivate broader community-reviewed prompt
construction across naturalistic AAVE, additional English varieties, and
multilingual settings.

The paired comparison prompts form a corpus condition. Their formality
varies across scenario areas, and the chest-pain triage prompts include
more clinical wording than the initial everyday/register prompts. This
variation is part of the archived audit surface and should be coded
explicitly in follow-up studies.

Current labels combine exact I/O review, heuristic prescreen columns,
and case evidence. A larger adjudicated study should add independent
coders, agreement statistics, bootstrap confidence intervals over
per-pair effects, and multiple-comparison procedures for aggregate
claims.

Deterministic decoding improves reproducibility and makes matched-pair
comparison cleaner. Sampled deployment behavior adds another layer:
temperature, nucleus sampling, system prompts, tool use, and moderation
layers may change refusal rates and divergence patterns. The
deterministic findings provide a reproducible baseline for those
extensions.

Visible traces are partial output surfaces. The dialect-to-location
claim is a claim about emitted trace content. Hidden causal pathways
remain open for future mechanistic work.

Hauhau is a refusal-reduced fine-tune with multiple behavioral
differences from the base model. The paper uses Hauhau to stress-test
response trajectories and to reveal how refusal convergence shapes
visible behavior.

\section{Ethics and Release}\label{ethics-and-release}

This work studies synthetic safety-test prompts involving self-harm,
violence, illegal activity, and slur-containing language. The dataset
does not contain personal user records or private user conversations.
The release concern is content amplification: public artifacts mask
slurs and quote only minimal evidence phrases when examples involve
operationally harmful completions. Exact prompts and outputs are
retained in a private controlled archive with hashes and provenance for
verification.

The purpose of DMRA is protective. Dialect-conditioned assistance
differences can disproportionately affect users already exposed to
unequal treatment in medical, legal, financial, and crisis-response
systems. Auditing these differences before deployment supports safer,
more equitable LLM assistance.

\section{Data and Code Availability}\label{data-and-code-availability}

The manuscript source and reproducibility metadata are prepared for
public release with the preprint. A sanitized public artifact bundle is
available at
\url{https://github.com/jeffreywilliamportfolio/dmra-public-artifacts}.
The bundle contains the run manifest, model identifiers and hashes,
deterministic decoding settings, checksums, and redacted case-evidence
summaries. The paper source, public metadata, and redacted public
summaries are released under CC BY 4.0. The public bundle excludes model
weights, environment files, secrets, full raw safety-test generations,
slur-containing prompt text, and operationally harmful completions.

Full raw generations are retained in a private controlled archive and
are available from the author upon reasonable request for verification.
The retained archive preserves file-level provenance and checksums while
keeping high-risk generated text separate from the public release.

\section{Conclusion}\label{conclusion}

Dialect-Marked Response Audit shows that matched safety-critical prompts
can receive different safety support when expressed through AAVE-marked
prompt surfaces and paired comparison prompts. The differences include
opposite epistemic framing, unequal medical scaffolding, different
empathy registers, visible dialect-to-location inference, and high-risk
tactical continuation under refusal-reduced settings.

The central claim is that safety evaluation needs the matched
interaction as its unit of analysis. Refusal decisions and correctness
labels capture important pieces of safety, while equivalent support
requires asking whether matched requests receive comparable urgency,
specificity, respect, and risk reduction. DMRA provides one protocol for
testing that equivalence. The observed divergences are concrete enough
to show that dialect equity in LLM safety requires evaluation of
safety-support equivalence across prompt surfaces.

\section*{Acknowledgments}\label{acknowledgments}
\addcontentsline{toc}{section}{Acknowledgments}

This work was self-funded. The author used OpenAI language-model tools
for drafting assistance, copyediting, LaTeX preparation, artifact
checks, and reference-scope review. The author reviewed the final
manuscript and takes responsibility for all claims, citations, analyses,
and text.

\section*{References}\label{references}
\addcontentsline{toc}{section}{References}

Blodgett, Su Lin, Lisa Green, and Brendan O'Connor. 2016. ``Demographic
Dialectal Variation in Social Media: A Case Study of African-American
English.'' In \emph{Proceedings of the 2016 Conference on Empirical
Methods in Natural Language Processing}, 1119-1130. Association for
Computational Linguistics. https://aclanthology.org/D16-1120/

Balasubramanian, Sriram, Samyadeep Basu, and Soheil Feizi. 2025. ``A
Closer Look at Bias and Chain-of-Thought Faithfulness of Large (Vision)
Language Models.'' arXiv:2505.23945. https://arxiv.org/abs/2505.23945

Dunlap, Deja, and R. Thomas McCoy. 2026. ``Evaluating the Usage of
African-American Vernacular English in Large Language Models.''
arXiv:2602.21485. https://arxiv.org/abs/2602.21485

Emmons, Scott, Roland S. Zimmermann, David K. Elson, and Rohin Shah.
2025. ``A Pragmatic Way to Measure Chain-of-Thought Monitorability.''
arXiv:2510.23966. https://arxiv.org/abs/2510.23966

Gehman, Samuel, Suchin Gururangan, Maarten Sap, Yejin Choi, and Noah A.
Smith. 2020. ``RealToxicityPrompts: Evaluating Neural Toxic Degeneration
in Language Models.'' In \emph{Findings of the Association for
Computational Linguistics: EMNLP 2020}, 3356-3369. Association for
Computational Linguistics.
https://aclanthology.org/2020.findings-emnlp.301/

Groenwold, Sophie, Lily Ou, Aesha Parekh, Samhita Honnavalli, Sharon
Levy, Diba Mirza, and William Yang Wang. 2020. ``Investigating
African-American Vernacular English in Transformer-Based Text
Generation.'' In \emph{Proceedings of the 2020 Conference on Empirical
Methods in Natural Language Processing}, 5877-5883. Association for
Computational Linguistics. https://aclanthology.org/2020.emnlp-main.473/

Gupta, Abhay, Jacob Cheung, Philip Meng, Shayan Sayyed, Austen Liao,
Kevin Zhu, and Sean O'Brien. 2025. ``EnDive: A Cross-Dialect Benchmark
for Fairness and Performance in Large Language Models.''
arXiv:2504.07100. https://arxiv.org/abs/2504.07100

Harvey, Emma, Rene F. Kizilcec, and Allison Koenecke. 2025. ``A
Framework for Auditing Chatbots for Dialect-Based Quality-of-Service
Harms.'' arXiv:2506.04419. https://arxiv.org/abs/2506.04419

Hofmann, Valentin, Pratyusha Ria Kalluri, Dan Jurafsky, and Sharese
King. 2024. ``AI Generates Covertly Racist Decisions About People Based
on Their Dialect.'' \emph{Nature} 633: 147-154.
https://doi.org/10.1038/s41586-024-07856-5

Khorramrouz, Adel, and Sharon Levy. 2025. ``Characterizing Selective
Refusal Bias in Large Language Models.'' arXiv:2510.27087.
https://arxiv.org/abs/2510.27087

Korbak, Tomek, Mikita Balesni, Elizabeth Barnes, Yoshua Bengio, Joe
Benton, Joseph Bloom, Mark Chen, et al.~2025. ``Chain of Thought
Monitorability: A New and Fragile Opportunity for AI Safety.''
arXiv:2507.11473. https://arxiv.org/abs/2507.11473

Kusner, Matt J., Joshua Loftus, Chris Russell, and Ricardo Silva. 2017.
``Counterfactual Fairness.'' In \emph{Advances in Neural Information
Processing Systems 30}.
https://papers.nips.cc/paper/6995-counterfactual-fairness

Lanham, Tamera, Anna Chen, Ansh Radhakrishnan, Benoit Steiner, Carson
Denison, Danny Hernandez, Dustin Li, Esin Durmus, Evan Hubinger, Jackson
Kernion, Kamile Lukosiute, Karina Nguyen, Newton Cheng, Nicholas Joseph,
Nicholas Schiefer, Oliver Rausch, Robin Larson, Sam McCandlish, Sandipan
Kundu, Saurav Kadavath, Shannon Yang, Thomas Henighan, Timothy Maxwell,
Timothy Telleen-Lawton, Tristan Hume, Zac Hatfield-Dodds, Jared Kaplan,
Jan Brauner, Samuel R. Bowman, and Ethan Perez. 2023. ``Measuring
Faithfulness in Chain-of-Thought Reasoning.'' arXiv:2307.13702.
https://arxiv.org/abs/2307.13702

Lin, Fangru, Shaoguang Mao, Emanuele La Malfa, Valentin Hofmann, Adrian
de Wynter, Xun Wang, Si-Qing Chen, Michael Wooldridge, Janet B.
Pierrehumbert, and Furu Wei. 2024. ``Assessing Dialect Fairness and
Robustness of Large Language Models in Reasoning Tasks.''
arXiv:2410.11005. https://arxiv.org/abs/2410.11005

Mazeika, Mantas, Long Phan, Xuwang Yin, Andy Zou, Zifan Wang, Norman Mu,
Elham Sakhaee, Nathaniel Li, Steven Basart, Bo Li, David Forsyth, and
Dan Hendrycks. 2024. ``HarmBench: A Standardized Evaluation Framework
for Automated Red Teaming and Robust Refusal.'' arXiv:2402.04249.
https://arxiv.org/abs/2402.04249

Maskey, Utsav, Mark Dras, and Usman Naseem. 2025. ``Should LLM Safety Be
More Than Refusing Harmful Instructions?'' arXiv:2506.02442.
https://arxiv.org/abs/2506.02442

Meek, Austin, Eitan Sprejer, Ivan Arcuschin, Austin J. Brockmeier, and
Steven Basart. 2025. ``Measuring Chain-of-Thought Monitorability Through
Faithfulness and Verbosity.'' arXiv:2510.27378.
https://arxiv.org/abs/2510.27378

Parrish, Alicia, Angelica Chen, Nikita Nangia, Vishakh Padmakumar, Jason
Phang, Jana Thompson, Phu Mon Htut, and Samuel R. Bowman. 2022. ``BBQ: A
Hand-Built Bias Benchmark for Question Answering.'' In \emph{Findings of
the Association for Computational Linguistics: ACL 2022}, 2086-2105.
Association for Computational Linguistics.
https://aclanthology.org/2022.findings-acl.165/

Qwen Team. 2025. ``Qwen3 Technical Report.'' arXiv:2505.09388.
https://arxiv.org/abs/2505.09388

Qwen Team. 2026. ``Qwen3.5: Towards Native Multimodal Agents.''
February. Model card for Qwen/Qwen3.5-35B-A3B, Hugging Face. Accessed
May 22, 2026. https://huggingface.co/Qwen/Qwen3.5-35B-A3B

Ribeiro, Marco Tulio, Tongshuang Wu, Carlos Guestrin, and Sameer Singh.
2020. ``Beyond Accuracy: Behavioral Testing of NLP Models with
CheckList.'' In \emph{Proceedings of the 58th Annual Meeting of the
Association for Computational Linguistics}, 4902-4912. Association for
Computational Linguistics. https://aclanthology.org/2020.acl-main.442/

Roettger, Paul, Hannah Rose Kirk, Bertie Vidgen, Giuseppe Attanasio,
Federico Bianchi, and Dirk Hovy. 2023. ``XSTest: A Test Suite for
Identifying Exaggerated Safety Behaviours in Large Language Models.''
arXiv:2308.01263. https://arxiv.org/abs/2308.01263

Sap, Maarten, Dallas Card, Saadia Gabriel, Yejin Choi, and Noah A.
Smith. 2019. ``The Risk of Racial Bias in Hate Speech Detection.'' In
\emph{Proceedings of the 57th Annual Meeting of the Association for
Computational Linguistics}, 1668-1678. Association for Computational
Linguistics. https://aclanthology.org/P19-1163/

Turpin, Miles, Julian Michael, Ethan Perez, and Samuel R. Bowman. 2023.
``Language Models Don't Always Say What They Think: Unfaithful
Explanations in Chain-of-Thought Prompting.'' In \emph{Advances in
Neural Information Processing Systems 36}.
https://arxiv.org/abs/2305.04388

Tornberg, Petter. 2026. ``Large Language Models Reproduce Racial
Stereotypes When Used for Text Annotation.'' arXiv:2603.13891.
https://arxiv.org/abs/2603.13891

Zou, Andy, Zifan Wang, Nicholas Carlini, Milad Nasr, J. Zico Kolter, and
Matt Fredrikson. 2023. ``Universal and Transferable Adversarial Attacks
on Aligned Language Models.'' arXiv:2307.15043.
https://arxiv.org/abs/2307.15043

\end{document}
